\problemname{Great Migration}

The swans are preparing for their great migration and need to maximise their energy reserves for the long flight.
Because they need to stay light enough to fly, a swan can only carry a maximum weight of food in their digestive tract.
The swans have a limited selection of aquatic plants available to eat.
Each plant has a specific weight and yields a specific amount of migratory energy.
A swan can choose to eat a plant or leave it, but they cannot eat a fraction of a plant.

\section*{Input}
The first line of the input consists of two integers representing the number of plants ($N$) and the total stomach capacity ($W$), bounded by ($1 \le N \le 1000$) and ($1 \le W \le 5000$).

Each of the next $N$ lines consists of two integers per line representing the weight ($w_i$) and energy value ($v_i$) of each plant, bounded by ($1 \le w_i \le W$) and ($1 \le v_i \le 10^4$).

\section*{Output}
Output a single integer representing the maximum total energy the swan can accumulate without exceeding its weight capacity.